\chapter{Marco europeo y lectura de normativa vigente}\label{ch:14}
\objetivos{Construir un mapa de los reglamentos principales y distinguir publicación, entrada en vigor, aplicación y consolidación.}
\section{Tres reglamentos, funciones diferentes}
El Reglamento 2018/1139 constituye la base común de seguridad aérea del sistema europeo. El Reglamento 748/2012 desarrolla, dentro de su ámbito, certificación o declaración de conformidad y requisitos de diseño y producción. El Reglamento 1321/2014 regula mantenimiento de la aeronavegabilidad y aprobación de organizaciones y personal involucrados. La separación facilita saber dónde buscar, aunque los documentos se remiten entre sí.

No es correcto tratar todas las aeronaves europeas como si cumplieran idéntico régimen. El artículo 2 del reglamento base delimita alcance, exclusiones y posibilidades específicas. Las aeronaves excluidas, las actividades estatales y la normativa nacional requieren un análisis propio. El marco militar se estudia en el capítulo~\ref{ch:23}.
\section{Versiones utilizadas en esta edición}
La consulta de esta edición se fija el \textbf{18 de octubre de 2026}. Se descargaron de EUR-Lex los textos consolidados 02018R1139--20260802, 02012R0748--20260222 y 02014R1321--20260807. La fecha del texto consolidado no es la fecha de aprobación del reglamento original.

Los textos consolidados facilitan la lectura, pero sus cabeceras indican que son instrumentos documentales sin efecto jurídico propio: las versiones auténticas son los actos del Diario Oficial. Cuando una modificación importa para una conclusión se consulta también el acto modificador, no solo una portada de consolidación.
\begin{table}[htbp]\centering\small\caption{Base documental normativa de la edición.}\label{tab:versiones}
\begin{tabular}{p{.2\linewidth}p{.28\linewidth}p{.4\linewidth}}\toprule
Norma & Consolidación consultada & Función de lectura\\\midrule
2018/1139 & 2 agosto 2026 & Ámbito, requisitos esenciales, sistema institucional\\[4pt]
748/2012 & 22 febrero 2026 & Parte 21 y certificación/declaraciones dentro de su ámbito\\[4pt]
1321/2014 & 7 agosto 2026 & Partes de aeronavegabilidad, mantenimiento y personal\\[4pt]
2026/100 & Acto publicado 19 enero 2026 & Cambios, entre otros, en revisión de aeronavegabilidad\\\bottomrule
\end{tabular}\end{table}
\section{Una modificación puede aplicarse después}
Un acto puede publicarse, entrar en vigor tras un plazo y contener fechas posteriores de aplicación o disposiciones transitorias. Es un error utilizar su fecha de publicación como si todas las disposiciones ya fueran exigibles ese día.

El Reglamento de Ejecución 2026/100, artículo 4, establece aplicación general desde el 7 de agosto de 2026 y fechas específicas para determinadas correcciones. Esta edición incorpora su efecto sobre la revisión de aeronavegabilidad al explicar M.A.901 y M.A.903. Las explicaciones de diapositivas antiguas sobre numeración o proceso no se trasplantan automáticamente al texto actual.
\section{Partes y reglas de aplicabilidad}
El Reglamento 1321/2014 contiene Parte M, Parte 145, Parte 66, Parte 147, Parte T, Parte ML, Parte CAMO y Parte CAO, entre otros elementos de su estructura. Cada parte tiene objeto y población de aplicación; el artículo 3 ayuda a elegir entre Parte M y Parte ML.

Por ejemplo, la categoría de aviones distintos de aeronaves motopropulsadas complejas con masa máxima de despegue de $2730\,\mathrm{kg}$ o menos aparece en el artículo 3.2, junto a otras categorías. El mismo artículo contiene una excepción para aeronaves incluidas en el AOC de una compañía aérea con licencia conforme al Reglamento 1008/2008. Por ello «menos de 2730 kg implica siempre Parte ML» es una simplificación incompleta.

La aprobación organizativa tampoco se obtiene solo de la masa. Se analizan operación, producto, tareas y régimen específico. Una tabla de resumen del libro orienta la búsqueda; el texto legal y los datos de la situación resuelven la conclusión.
\section{AMC, GM, CS y Easy Access Rules}
Los Easy Access Rules de EASA reúnen material para una lectura práctica, incluyendo reglas, AMC y GM. Debe comprobarse su fecha y sus avisos de actualización. La revisión de aeronavegabilidad consultada de septiembre de 2025 declara que no incorpora el Reglamento 2025/111; no debe presentarse como una recopilación completa de toda norma aplicable en octubre de 2026.

La página oficial de CS-E muestra enmienda 9 publicada el 25 de septiembre de 2026. El PDF está fechado el 21 de septiembre y remite a la decisión correspondiente para su entrada en vigor. Esto demuestra otra distinción: fecha de documento, publicación y aplicabilidad no son sinónimos. La base de certificación de un motor específico puede conservar una enmienda distinta según el proceso aplicable.
\caso{M14. La copia antigua no decide la vigencia}{Se recibe un PDF de 1321/2014 cuyo encabezado muestra una consolidación anterior y una explicación de ARC de 2022. Se conserva como fuente histórica/docente, se identifica la versión actual y se compara el apartado pertinente. La conclusión se apoya en M.A.901 y M.A.903 de la consolidación de agosto de 2026 y en el acto modificador cuando procede. No se destruye la copia antigua: se etiqueta su alcance para evitar confundirla con la referencia vigente.}
\section{Cómo registrar una afirmación jurídica}
La ficha mínima contiene autoridad y jurisdicción, norma, versión, punto exacto, tipo de obligación, fecha de aplicación, fuente oficial y conclusión. Una matriz separa afirmación contrastada, explicación académica y cuestión pendiente. No se declara «normativa verificada» por haber descargado un reglamento sin comprobar los artículos usados.

El informe normativo de la entrega enumera las disposiciones de las que dependen las afirmaciones importantes. Las modificaciones futuras, la configuración de una aeronave real y los criterios de una autoridad para un expediente concreto quedan fuera de una edición cerrada en esta fecha.
\section{Materiales normativos históricos y alcance de sus conclusiones}
Los documentos aportados contienen consolidaciones de 1321/2014 fechadas en 2019 y de 2018/1139 fechadas en 2021. Son útiles para seguir la evolución, pero no bastan para responder sobre el derecho aplicable en octubre de 2026. Se contrasta la versión pertinente, el artículo de ámbito, la modificación y su fecha de aplicación. Deben distinguirse ley, reglamento, AMC, GM y especificaciones de certificación.

En un supuesto de examen se determina primero ámbito civil o estatal, clase y operación de aeronave, Parte M o ML cuando corresponda, responsables, documento y momento relevante. No se atribuyen idénticos privilegios a Parte 21, Parte 145, Parte 66, Parte 147, CAMO y CAO.


\section{Citar de forma reproducible}
«Según EASA» no permite reproducir una conclusión. La ficha de esta obra usa autoridad, acto, versión, punto, ámbito, fecha y observaciones. La conclusión añade hechos que activan el requisito: masa, complejidad, operación y configuración, según corresponda. Si falta un hecho decisivo se conserva la condición pendiente.

Una consolidación facilita lectura; el acto modificativo ayuda a identificar qué cambió y desde cuándo. Entrada en vigor y aplicación pueden diferir y contener excepciones. Se comprueba la cláusula en lugar de inferirla del día de publicación o del año del reglamento. También se separan legislación y material explicativo.
\section{La modificación de 2026 como ejercicio de control de versiones}
El Reglamento de ejecución 2026/100 modifica, entre otros aspectos, el marco de revisión de aeronavegabilidad. Su artículo 4 fija aplicación general desde el 7 de agosto de 2026 con excepciones específicas. M.A.901 y M.A.903 se leen en la redacción actual. Una tabla de 2025 puede mantener conceptos útiles y describir una numeración o proceso que exige actualización.

Esta edición registra consolidación 1321/2014 de 7 de agosto de 2026 y acto 2026/100. La EAR de septiembre de 2025 se usa con su fecha, sin atribuirle todos los cambios posteriores. La discrepancia se resuelve por disposición, sin declarar inútil todo documento antiguo.
\section{Matriz de aplicación}
Para escoger M o ML se documentan las condiciones del artículo 3; «avión pequeño» es insuficiente. La excepción relativa a aeronaves incluidas en el AOC de una compañía aérea con licencia 1008/2008 obliga a combinar características y operación. Después se analiza qué parte regula programa, gestión, ejecución y revisión; no se trasladan todas las conclusiones de M a ML.

Una norma general tampoco contiene los límites técnicos de cada motor. Si falta ese dato, se mantiene un valor ficticio claramente identificado, sin presentarlo como límite certificado.
\repaso{¿La norma descargada hoy prueba cumplimiento en un hecho del año pasado?}{Se reconstruyen versión y requisitos aplicables entonces y evidencia de ejecución. El texto actual responde a otra fecha y puede haber cambiado.}

\resumen{Primero se determina el ámbito; después se elige la norma y su versión. Consolidación no equivale a acto auténtico, publicación no equivale a aplicación y última enmienda no equivale a base de certificación de todos los productos. La documentación antigua se conserva con su fecha y función.}
\fuentes{EUR-Lex \cite{EU1139,EU748,EU1321,EU2026-100}; EASA, aviso de EAR de septiembre de 2025 y CS-E enmienda 9 \cite{EASA-CAW,CSE9}.}
